%% Springer Nature journal template (sn-jnl), numeric math/physics reference style.
%% Requires sn-jnl.cls and sn-mathphys-num.bst from the official Springer Nature LaTeX template
%% (https://www.springernature.com/gp/authors/campaigns/latex-author-support). Compile:
%%   pdflatex sn-checkpoint1 && bibtex sn-checkpoint1 && pdflatex sn-checkpoint1 && pdflatex sn-checkpoint1
\documentclass[pdflatex,sn-mathphys-num]{sn-jnl}
\usepackage{graphicx,multirow,amsmath,amssymb,amsfonts,amsthm,mathrsfs,xcolor,textcomp,manyfoot,booktabs}
\usepackage{longtable,array,calc,pifont}
\usepackage[title]{appendix}
\providecommand{\tightlist}{\setlength{\itemsep}{0pt}\setlength{\parskip}{0pt}}
\raggedbottom

\begin{document}

\title[Checkpoint 1: Slice-Aware Risk-Controlled Intrusion Response]{Discovering a Defensible 5G/6G Research Problem: Checkpoint 1 Report on Certified Collateral-Risk Control for Slice-Aware Intrusion Response}

\author*[1]{\fnm{[Author]} \sur{[To be supplied]}}\email{[to be supplied]}
\affil*[1]{\orgname{[Affiliation to be supplied]}}

\abstract{We report Checkpoint 1 of an independent research-discovery mission whose goal is a novel, feasible, and defensible IEEE Transactions-level problem in 5G/6G networks. We screened 75 sources across ten research clusters; 72 papers and preprints are usable, each recorded with verification status and provenance in a reusable structured bibliography. From them we generated eleven candidate problems, scored them with a fixed weighted rubric, and applied mandatory rejection criteria. Six candidates were rejected and five retained; the top three went to an adversarial novelty audit. The audit shows that detection-level conformal false-alarm control for intrusion detection, the direction closest to the team's prior work, became crowded in 2025-2026, so we reject it as a primary topic. We provisionally recommend certified collateral-risk control for slice-aware intrusion response: a multi-action mitigation policy calibrated, through multiple-testing risk control, so that the oracle-counterfactual SLA-violation risk imposed on benign users of each slice stays within a slice-specific budget with a finite-sample, distribution-free guarantee. Our backup is SLA-weighted multi-slice quickest attack detection. The report states falsifiable hypotheses, baselines, and minimum evidence, and lists the access gaps that block final novelty claims. It contains no experimental results.}

\keywords{5G network slicing, intrusion response, conformal risk control, learn-then-test, service-level agreements, research methodology}

\maketitle

\hypertarget{sec:intro}{%
\section{Introduction}\label{sec:intro}}

This report is \textbf{Checkpoint 1} of an independent
research-discovery mission. The mission asks for a novel, feasible, and
defensible IEEE Transactions-level research problem in 5G/6G systems.
Its deliverable is a \emph{validated research direction and an
executable plan}, not a manuscript with results. As the brief requires,
the report stops after the initial checkpoint (Stages A--F, preliminary
forms only). It contains \textbf{no experimental results and no
performance numbers} for any proposed method. Every claim is tagged with
one of four evidence levels:

\begin{itemize}
\tightlist
\item
  \textbf{{[}Verified{]}}: confirmed in a retrieved bibliographic or
  repository record.
\item
  \textbf{{[}Inference{]}}: an evidence-supported inference by the
  authors of this report.
\item
  \textbf{{[}Hypothesis{]}}: a falsifiable statement that has not yet
  been tested.
\item
  \textbf{{[}Unverified{]}}: could not be checked in this session.
\end{itemize}

\textbf{Summary of the provisional outcome.} We screened 75 sources
across ten 5G/6G research areas. Of these, 72 papers and preprints are
usable; 3 entries are excluded from that count (2 software artefacts and
1 unconfirmed preprint). From them we generated eleven candidate
problems. Six were rejected under the mandatory criteria and five
retained; the three highest-scoring retained candidates (C1, C3, C6)
went to an adversarial novelty audit. The most important negative
finding is this: \emph{detection-level} conformal false-alarm control
for intrusion detection, the direction that best matched the team's
prior background, has become crowded in 2025--2026
\citep{youssef2026caliburn, barrett2026firce, confid2026, robustcids2026, farzaneh2025ccoad}.
It is therefore \textbf{rejected} as a primary topic.

The provisional primary recommendation moves the guarantee one layer
down the control loop, from \emph{alerts} to the \emph{service-level
consequences of mitigation actions}:

\begin{quote}
\textbf{C1 --- Risk-controlled, slice-aware intrusion response.}
Calibrate detector-triggered mitigation policies (rate-limit, quarantine
to a sinkhole slice, block) for a sliced 5G core. The calibration should
come with a finite-sample, distribution-free guarantee that the
\emph{collateral SLA-violation risk imposed on benign users of each
slice} stays below a slice-specific budget, while residual attack impact
is minimised.
\end{quote}

The backup is \textbf{C3}, SLA-weighted multi-slice quickest attack
detection. Both recommendations are explicitly \textbf{provisional}.
Section \protect\hyperlink{sec:gaps}{9} lists the full-text reads and
infrastructure checks that must pass before Phase 2.

\hypertarget{sec:method}{%
\section{Methodology of this checkpoint}\label{sec:method}}

\textbf{Search protocol.} We ran targeted web searches (about 60
queries, standard and extended) on 2 November 2026. They covered ten
topic clusters from the brief (Section
\protect\hyperlink{sec:landscape}{3}), plus explicit adversarial
searches for each shortlisted claim (Section
\protect\hyperlink{sec:audit}{6}). For each source we recorded the
citation, venue, year, DOI or preprint identifier, the problem,
contribution, method, environment, findings, limitations, relevance, and
verification status. The structured bibliography is
\texttt{data/bibliography.json}. The matrix
(\texttt{data/literature\_matrix.csv}) and the BibTeX file are generated
from it by \texttt{scripts/build\_bib.py}.

\textbf{Verification levels.} Each entry is labelled V-DOI, V-META,
V-PRE, PARTIAL, or UNVERIFIED (Table 1). \textbf{No DOI was guessed.}
Where a DOI did not appear in any retrieved record, the field is left
blank.

\textbf{Hard constraints of this session {[}Verified{]}.} The
container's network policy blocked arxiv.org, doi.org, api.crossref.org,
OpenAlex, and Semantic Scholar (HTTP 403 or DNS failure). Only a
search-engine interface was available. As a result:

\begin{enumerate}
\def\labelenumi{\arabic{enumi}.}
\tightlist
\item
  \textbf{No full text was read.} Every entry was reviewed at the
  abstract and metadata level. Limitations marked {[}A{]} in the matrix
  are the authors' own statements; those marked {[}I{]} are our
  inferences and must be confirmed against the full text.
\item
  DOIs could not be resolved independently.
\item
  The repository supplied for this mission was \textbf{empty}: no prior
  manuscripts, code, configurations, logs, or datasets. Claims about the
  team's prior work and testbed therefore stay
  \textbf{{[}Unverified{]}}.
\end{enumerate}

\begin{longtable}[]{@{}
  >{\raggedright\arraybackslash}p{(\columnwidth - 4\tabcolsep) * \real{0.3333}}
  >{\raggedright\arraybackslash}p{(\columnwidth - 4\tabcolsep) * \real{0.3333}}
  >{\raggedleft\arraybackslash}p{(\columnwidth - 4\tabcolsep) * \real{0.3333}}@{}}
\caption{Verification status of the structured
bibliography.}\tabularnewline
\toprule\noalign{}
\begin{minipage}[b]{\linewidth}\raggedright
Level
\end{minipage} & \begin{minipage}[b]{\linewidth}\raggedright
Meaning
\end{minipage} & \begin{minipage}[b]{\linewidth}\raggedleft
Count
\end{minipage} \\
\midrule\noalign{}
\endfirsthead
\toprule\noalign{}
\begin{minipage}[b]{\linewidth}\raggedright
Level
\end{minipage} & \begin{minipage}[b]{\linewidth}\raggedright
Meaning
\end{minipage} & \begin{minipage}[b]{\linewidth}\raggedleft
Count
\end{minipage} \\
\midrule\noalign{}
\endhead
\bottomrule\noalign{}
\endlastfoot
V-DOI & Title, authors, venue, year, and DOI seen together in a
retrieved record & 24 \\
V-META & Metadata verified; DOI not seen (left blank) & 20 \\
V-PRE & Preprint (or software artefact) exists; no peer-reviewed version
confirmed & 19 \\
PARTIAL & Exists; at least one metadata field conflicts or is
unconfirmed & 11 \\
UNVERIFIED & Seen once; must not be cited until checked & 1 \\
\textbf{Total} & 45 of the 75 entries are dated 2023--2026; 7 are
surveys used only as background & \textbf{75} \\
\end{longtable}

\textbf{Honest assessment of the minimum literature target.} The brief
asks for at least 60 \emph{verified} relevant papers. We have 72 usable
papers and preprints, so the count is met. The \emph{depth} of
verification is weaker than a Transactions-grade review requires: only
24 entries have a confirmed DOI, 19 are preprints, and none was read in
full. We therefore treat the landscape as \textbf{sufficient for
screening and ranking, but not for final novelty claims}. Phase 1 of the
roadmap (Section \protect\hyperlink{sec:roadmap}{8}) closes this gap.

\hypertarget{sec:landscape}{%
\section{Stage A --- Research landscape}\label{sec:landscape}}

Table 2 organises the evidence into ten clusters. The full per-paper
matrix (75 rows $\times$ 19 fields) is supplied as
\texttt{data/literature\_matrix.csv}.

\begin{longtable}[]{@{}
  >{\centering\arraybackslash}p{(\columnwidth - 8\tabcolsep) * \real{0.2000}}
  >{\raggedright\arraybackslash}p{(\columnwidth - 8\tabcolsep) * \real{0.2000}}
  >{\raggedright\arraybackslash}p{(\columnwidth - 8\tabcolsep) * \real{0.2000}}
  >{\raggedright\arraybackslash}p{(\columnwidth - 8\tabcolsep) * \real{0.2000}}
  >{\raggedright\arraybackslash}p{(\columnwidth - 8\tabcolsep) * \real{0.2000}}@{}}
\caption{Structured research landscape.}\tabularnewline
\toprule\noalign{}
\begin{minipage}[b]{\linewidth}\centering
\#
\end{minipage} & \begin{minipage}[b]{\linewidth}\raggedright
Cluster
\end{minipage} & \begin{minipage}[b]{\linewidth}\raggedright
Representative verified works
\end{minipage} & \begin{minipage}[b]{\linewidth}\raggedright
State of the art (evidence level)
\end{minipage} & \begin{minipage}[b]{\linewidth}\raggedright
Open space
\end{minipage} \\
\midrule\noalign{}
\endfirsthead
\toprule\noalign{}
\begin{minipage}[b]{\linewidth}\centering
\#
\end{minipage} & \begin{minipage}[b]{\linewidth}\raggedright
Cluster
\end{minipage} & \begin{minipage}[b]{\linewidth}\raggedright
Representative verified works
\end{minipage} & \begin{minipage}[b]{\linewidth}\raggedright
State of the art (evidence level)
\end{minipage} & \begin{minipage}[b]{\linewidth}\raggedright
Open space
\end{minipage} \\
\midrule\noalign{}
\endhead
\bottomrule\noalign{}
\endlastfoot
L1 & Validity of ML-NIDS evaluation &
\citep{sommer2010outside, arp2022dos, pendlebury2019tesseract, engelen2021troubleshooting, flood2024smells, apruzzese2023sok, jacobs2022emperor}
& Benchmark results are inflated by sampling bias, temporal leakage,
label errors, and shortcut features. There is a broad consensus that
false-positive cost, not accuracy, blocks deployment. {[}Verified,
abstract level{]} & Operational, cost-aware evaluation is still rare. \\
L2 & Cross-domain generalisation &
\citep{layeghy2023xcross, sarhan2022standard, farzaneh2024dtl5g} &
Detectors degrade sharply across datasets; NetFlow feature
standardisation enables cross-testing. {[}Verified{]} & No 5G
slice-aware cross-domain study found. {[}Inference{]} \\
L3 & 5G-specific security data and detection &
\citep{samarakoon2022nidd, siriwardhana2025nidd, khan2023slicesecure, amponis2022pfcp, manca2026sage5gc, slicespoof2025, xu2025rogue}
& Detection dominates. Reported accuracies are near-perfect (for example
99.99\% in \citep{khan2023slicesecure}, per its abstract). SAGE-5GC
shows adversarial fragility. {[}Verified{]} & \textbf{Response and
mitigation, and their service-level side effects, are almost absent.}
{[}Inference{]} \\
L4 & Drift and conformal security &
\citep{yang2021cade, han2023owad, barbero2022transcendent, barrett2026firce, youssef2026caliburn, robustcids2026, confid2026, farzaneh2025ccoad, rbsid2025}
& 2025--2026 saw a surge of conformal and risk-calibrated IDS: FPR or
FDR control, regime-dependent CRC, robustness to adversarial orbits.
{[}Verified, preprint level{]} & Detection-level guarantees are now
crowded; \emph{action-level} guarantees are not. {[}Inference{]} \\
L5 & Distribution-free risk-control theory &
\citep{tibshirani2019weighted, barber2023beyond, gibbs2021aci, gibbs2024online, bates2023outliers, angelopoulos2024crc, angelopoulos2025ltt, crcnonmono2026, cora2026, cpc2026}
& Mature tools exist for monotone (CRC) and non-monotone or multiple
(LTT) risks, covariate shift, and online shift. Action gating with risk
control is appearing in other domains (GUI agents). {[}Verified{]} & Not
yet applied to coupled, multi-action network mitigation, as far as
found. {[}Inference{]} \\
L6 & Intrusion response and mitigation &
\citep{hammar2022optimalstopping, hammar2024dynamic, hammar2023recursive, hammar2025conjectural, zonouz2014rre, regard2023, liu2021jaqen}
& Model-based optimal response (POMDPs, stochastic games, optimal
stopping) yields threshold policies; line-rate DDoS mitigation exists in
P4. {[}Verified{]} & No distribution-free guarantee on \emph{collateral}
harm; no slice awareness. {[}Inference{]} \\
L7 & Slicing and SLA-aware resource management &
\citep{foukas2017survey, foukas2017orion, afolabi2018slicing, bega2019deepcog, salvat2018overbooking, proactiveassurance2024}
& Asymmetric SLA-violation costs and overbooking create cross-slice
coupling. {[}Verified{]} & Coupling between security actions and SLAs is
unexplored. {[}Inference{]} \\
L8 & O-RAN and AI-native RAN control &
\citep{polese2023understanding, polese2023coloran, bonati2023openrangym, bonati2021colosseum, fiandrino2023explora, brachdelprever2025pacifista, groen2024implementing, groen2024securing}
& xApp development, explainability, conflict management (PACIFISTA, TMC
2025), interface security. {[}Verified{]} & All require RAN or RIC
infrastructure that is not available here. \\
L9 & Adversarial attacks on learned slicing &
\citep{shi2021attackdefend, shi2021flooding, tashman2026adversarial, wang2023robustslicing}
& Jamming and flooding against DRL slicing cause slice-dependent SLA
violations. {[}Verified, preprint level{]} & Crowded, and it needs PHY
access. \\
L10 & Digital twins, energy, FL, LLM, QCD &
\citep{ferriolgalmes2023routenetfermi, almasan2022ndt, lopezperez2022energy, lavaur2022fl, flaccshield2026, zhou2024llmtelecom, boateng2025llm, tartakovsky2006novel, fellouris2016second, qcdslice_github}
& GNN twins give point predictions. Energy work needs RAN measurements.
FL-IDS and LLM-ops are crowded or hype-driven. Multistream QCD theory is
mature. {[}Verified{]} & Calibrated twins (C6) and SLA-weighted
multi-slice QCD (C3) are moderately open. {[}Inference{]} \\
\end{longtable}

\textbf{Cross-cutting observations.}

\begin{enumerate}
\def\labelenumi{\arabic{enumi}.}
\tightlist
\item
  \textbf{{[}Inference{]}} The 5G security literature is stuck at the
  \emph{detection} stage on near-saturated benchmarks. Clusters L1 and
  L4 already explain why accuracy claims are unreliable.
\item
  \textbf{{[}Verified{]}} Mature distribution-free risk-control tools
  (L5) exist, and conformal IDS work (L4) uses them, but only on
  \emph{alert-level} quantities such as FPR, FDR, and alert budgets.
\item
  \textbf{{[}Inference{]}} The \emph{service-level consequence} of
  acting on alerts is what operators pay for in a sliced network, and
  especially in URLLC. It is governed by model-based response theory
  (L6) that offers no finite-sample guarantee under model
  misspecification. It is also coupled across slices through shared
  resources (L7).
\end{enumerate}

The intersection of points 2 and 3 is the gap pursued by C1.

\textbf{Coverage limitations.} PHY-layer and IEEE TWC-style original
research (beamforming, ISAC, RIS) is under-represented. It was screened
out early because nothing indicates that the team has PHY or SDR
infrastructure (Section \protect\hyperlink{sec:feas}{7}). This is a
deliberate scoping decision, recorded in the decision log as D-04, not
evidence that PHY topics are less important.

\hypertarget{sec:candidates}{%
\section{Stage B --- Candidate research problems}\label{sec:candidates}}

We generated eleven candidates, one more than the ten required, so that
the rejected C11 could be documented. Each is summarised along the
brief's five dimensions: A (problem), B (novelty), C (technical depth),
D (feasibility), and E (evaluation). The shortlisted candidates are
expanded in Sections \protect\hyperlink{sec:audit}{6} and
\protect\hyperlink{sec:rec}{8}.

\textbf{C1 --- Risk-controlled, slice-aware intrusion response.}

\begin{itemize}
\tightlist
\item
  \emph{A.} Mitigation actions taken on imperfect detections
  (rate-limiting, quarantine, blocking) harm benign users. In a sliced
  core this harm propagates through shared UPF and transport resources,
  and its cost differs sharply by slice (URLLC versus mMTC). Current IDS
  work bounds alert-level error but not the SLA damage of the actions
  taken.
\item
  \emph{B.} Closest works: CALIBURN \citep{youssef2026caliburn}
  (FP-budgeted alerts), FIRCE \citep{barrett2026firce} (``response''
  means retraining), CORA \citep{cora2026} (risk-controlled action
  gating, non-network), Hammar \& Stadler
  \citep{hammar2022optimalstopping} (model-based optimal response), and
  SliceSecure \citep{khan2023slicesecure} (slice detection only).
  \textbf{Missing piece {[}Inference{]}:} a distribution-free, per-slice
  guarantee on the \emph{collateral SLA risk of a multi-action
  mitigation policy} under cross-slice coupling.
\item
  \emph{C.} Multi-risk LTT calibration over a structured policy space,
  counterfactual loss evaluation with a twin-error propagation bound,
  and robustness to temporal dependence and shift.
\item
  \emph{D.} Needs an emulated sliced 5G core (Open5GS or free5GC with
  UERANSIM), a traffic and attack harness, and KPI measurement. Moderate
  to high complexity. Main risk: the cost of counterfactual replay.
\item
  \emph{E.} See Section \protect\hyperlink{sec:rec}{8}.
\end{itemize}

\textbf{C2 --- Slice-stratified conformal false-alarm control for 5G
NIDS (detection only).}

\begin{itemize}
\tightlist
\item
  \emph{A.} Bound the per-slice FPR of a detector.
\item
  \emph{B.} \textbf{Contradicted.} Marginal FPR and FDR control
  \citep{bates2023outliers}, conformal IDS
  \citep{confid2026, barbero2022transcendent, barrett2026firce},
  regime-dependent CRC \citep{youssef2026caliburn}, and online conformal
  FDR \citep{farzaneh2025ccoad} already exist. Stratifying by slice is a
  Mondrian-style re-grouping and too thin for a Transactions paper.
\item
  \emph{C.} Low.
\item
  \emph{D.} High feasibility.
\item
  \emph{E.} Straightforward.
\item
  \textbf{Rejected:} gap contradicted.
\end{itemize}

\textbf{C3 --- SLA-weighted multi-slice quickest attack detection.}

\begin{itemize}
\tightlist
\item
  \emph{A.} Detect attack onset across slice streams so as to minimise
  \emph{SLA-weighted} detection delay subject to \emph{heterogeneous
  per-slice} false-alarm constraints, with streams coupled by shared
  resources.
\item
  \emph{B.} Closest works:
  \citep{tartakovsky2006novel, fellouris2016second}, plus an unpublished
  simulation \citep{qcdslice_github}. \textbf{Missing piece
  {[}Inference{]}:} heterogeneous delay weights and false-alarm
  constraints with coupled streams.
\item
  \emph{C.} Asymptotic optimality analysis is possible and necessary.
\item
  \emph{D.} Trace-driven evaluation is feasible.
\item
  \emph{E.} Delay versus ARL curves against CUSUM and SUM-CUSUM
  variants.
\item
  \textbf{Risk:} the result may be a corollary of the existing
  multichannel theory.
\end{itemize}

\textbf{C4 --- Leakage-audited cross-testbed benchmark for 5G NIDS.}

\begin{itemize}
\tightlist
\item
  \textbf{Rejected.} It probably overlaps the team's own leakage-audit
  work {[}Unverified: prior work not provided{]}, and a 5G-NIDD
  realistic-split package already exists \citep{tvquynh2026repro}.
  Valuable as \emph{infrastructure} for C1, not as the contribution.
\end{itemize}

\textbf{C5 --- Conflict-aware xApp composition with guarantees.}

\begin{itemize}
\tightlist
\item
  \textbf{Rejected.} PACIFISTA \citep{brachdelprever2025pacifista}
  covers conflict detection and severity, and no RAN or RIC
  infrastructure has been verified.
\end{itemize}

\textbf{C6 --- Calibrated-uncertainty GNN digital twin for slice SLA
prediction.}

\begin{itemize}
\tightlist
\item
  \emph{A.} GNN twins \citep{ferriolgalmes2023routenetfermi} give point
  estimates; operators need calibrated SLA-violation risk.
\item
  \emph{B.} No conformal-twin work was found in this session's searches,
  but \textbf{the search was shallow {[}Unverified{]}}.
\item
  \emph{C.} Conformal regression over graph-structured,
  exchangeability-breaking data.
\item
  \emph{D.} Feasible on public RouteNet datasets, but the 5G relevance
  is indirect (wired topologies).
\item
  \emph{E.} Coverage and width against Bayesian and ensemble GNNs.
\end{itemize}

\textbf{C7 --- Adversarial robustness of DRL RAN-slicing controllers.}

\begin{itemize}
\tightlist
\item
  \textbf{Rejected.} It needs PHY, RAN, or SDR access, and the topic is
  crowded
  \citep{shi2021attackdefend, tashman2026adversarial, wang2023robustslicing}.
\end{itemize}

\textbf{C8 --- Uncertainty-aware RAN sleep-mode control.}

\begin{itemize}
\tightlist
\item
  \textbf{Rejected.} It needs RAN energy measurements
  \citep{lopezperez2022energy}; in simulation alone it cannot support
  operational claims.
\end{itemize}

\textbf{C9 --- Federated slice-level IDS under non-IID data.}

\begin{itemize}
\tightlist
\item
  \textbf{Rejected.} Crowded \citep{lavaur2022fl, flaccshield2026}, with
  no identifiable new contribution.
\end{itemize}

\textbf{C10 --- Adversarially robust 5GC PFCP anomaly detection.}

\begin{itemize}
\tightlist
\item
  \emph{Retained but not shortlisted.} SAGE-5GC \citep{manca2026sage5gc}
  has just defined the evaluation protocol; a defence paper would be
  incremental unless it adds certified robustness
  \citep{robustcids2026}.
\end{itemize}

\textbf{C11 --- LLM-driven intent-based slice orchestration.}

\begin{itemize}
\tightlist
\item
  \textbf{Rejected.} Hype-driven, with no falsifiable evaluation
  possible without operator intents and data
  \citep{zhou2024llmtelecom, boateng2025llm}.
\end{itemize}

\hypertarget{sec:rank}{%
\section{Stage C --- Weighted ranking}\label{sec:rank}}

Scores run from 1 to 10 and use the brief's weights: novelty 25\%,
significance 20\%, feasibility 20\%, rigour 15\%, validation 10\%, fit
10\%. They are computed reproducibly by \texttt{scripts/score.py} from
\texttt{data/candidate\_scores.csv}. \textbf{The weighted total
summarises our judgement; it is not evidence of publishability.}

\begin{longtable}[]{@{}
  >{\centering\arraybackslash}p{(\columnwidth - 18\tabcolsep) * \real{0.1000}}
  >{\centering\arraybackslash}p{(\columnwidth - 18\tabcolsep) * \real{0.1000}}
  >{\centering\arraybackslash}p{(\columnwidth - 18\tabcolsep) * \real{0.1000}}
  >{\centering\arraybackslash}p{(\columnwidth - 18\tabcolsep) * \real{0.1000}}
  >{\centering\arraybackslash}p{(\columnwidth - 18\tabcolsep) * \real{0.1000}}
  >{\centering\arraybackslash}p{(\columnwidth - 18\tabcolsep) * \real{0.1000}}
  >{\centering\arraybackslash}p{(\columnwidth - 18\tabcolsep) * \real{0.1000}}
  >{\centering\arraybackslash}p{(\columnwidth - 18\tabcolsep) * \real{0.1000}}
  >{\centering\arraybackslash}p{(\columnwidth - 18\tabcolsep) * \real{0.1000}}
  >{\raggedright\arraybackslash}p{(\columnwidth - 18\tabcolsep) * \real{0.1000}}@{}}
\caption{Weighted ranking. Hard-rejected candidates are placed below
retained ones regardless of score.}\tabularnewline
\toprule\noalign{}
\begin{minipage}[b]{\linewidth}\centering
Rank
\end{minipage} & \begin{minipage}[b]{\linewidth}\centering
ID
\end{minipage} & \begin{minipage}[b]{\linewidth}\centering
Nov.
\end{minipage} & \begin{minipage}[b]{\linewidth}\centering
Sig.
\end{minipage} & \begin{minipage}[b]{\linewidth}\centering
Feas.
\end{minipage} & \begin{minipage}[b]{\linewidth}\centering
Rig.
\end{minipage} & \begin{minipage}[b]{\linewidth}\centering
Val.
\end{minipage} & \begin{minipage}[b]{\linewidth}\centering
Fit
\end{minipage} & \begin{minipage}[b]{\linewidth}\centering
Weighted
\end{minipage} & \begin{minipage}[b]{\linewidth}\raggedright
Decision
\end{minipage} \\
\midrule\noalign{}
\endfirsthead
\toprule\noalign{}
\begin{minipage}[b]{\linewidth}\centering
Rank
\end{minipage} & \begin{minipage}[b]{\linewidth}\centering
ID
\end{minipage} & \begin{minipage}[b]{\linewidth}\centering
Nov.
\end{minipage} & \begin{minipage}[b]{\linewidth}\centering
Sig.
\end{minipage} & \begin{minipage}[b]{\linewidth}\centering
Feas.
\end{minipage} & \begin{minipage}[b]{\linewidth}\centering
Rig.
\end{minipage} & \begin{minipage}[b]{\linewidth}\centering
Val.
\end{minipage} & \begin{minipage}[b]{\linewidth}\centering
Fit
\end{minipage} & \begin{minipage}[b]{\linewidth}\centering
Weighted
\end{minipage} & \begin{minipage}[b]{\linewidth}\raggedright
Decision
\end{minipage} \\
\midrule\noalign{}
\endhead
\bottomrule\noalign{}
\endlastfoot
1 & C1 & 7 & 8 & 7 & 8 & 7 & 8 & 7.45 & Shortlist (primary) \\
2 & C3 & 6 & 7 & 7 & 8 & 6 & 7 & 6.80 & Shortlist (backup) \\
3 & C6 & 6 & 7 & 6 & 7 & 6 & 7 & 6.45 & Shortlist \\
4 & C10 & 4 & 7 & 6 & 6 & 6 & 6 & 5.70 & Retained \\
5 & C4 & 4 & 6 & 8 & 5 & 7 & 5 & 5.75 & Rejected (overlap) \\
6 & C2 & 3 & 6 & 8 & 6 & 7 & 5 & 5.65 & Rejected (gap contradicted) \\
7 & C8 & 5 & 8 & 2 & 6 & 3 & 7 & 5.15 & Rejected (infrastructure) \\
8 & C5 & 4 & 8 & 2 & 7 & 3 & 7 & 5.05 & Rejected (prior work and
infrastructure) \\
9 & C9 & 3 & 5 & 7 & 5 & 5 & 4 & 4.80 & Rejected (no contribution) \\
10 & C7 & 4 & 7 & 2 & 6 & 3 & 6 & 4.60 & Rejected (infrastructure and
crowding) \\
11 & C11 & 3 & 6 & 5 & 4 & 4 & 5 & 4.45 & Rejected (unfalsifiable) \\
\end{longtable}

\textbf{Score justifications for the shortlist.}

\begin{itemize}
\tightlist
\item
  \textbf{C1, novelty 7 rather than 9.} The statistical mechanism (LTT)
  and the idea of risk-controlled action gating \citep{cora2026} both
  exist. Novelty must come from the \emph{network-specific structure}:
  per-slice budgets, multi-action policies, cross-slice coupling, and
  counterfactual evaluation. That novelty is plausible but has not been
  confirmed against full texts of
  \citep{youssef2026caliburn, rbsid2025, cpc2026}.
\item
  \textbf{C1, feasibility 7.} An emulated core is standard practice, but
  the existing testbed has not been audited.
\item
  \textbf{C3, rigour 8, novelty 6.} Strong theory is possible, but
  \citep{fellouris2016second} may already subsume the core result.
\item
  \textbf{C6, feasibility 6.} Public datasets exist, but the 5G link is
  indirect.
\end{itemize}

\hypertarget{sec:audit}{%
\section{Stage D --- Adversarial novelty audit}\label{sec:audit}}

\hypertarget{c1-risk-controlled-slice-aware-intrusion-response}{%
\subsection{C1: risk-controlled slice-aware intrusion
response}\label{c1-risk-controlled-slice-aware-intrusion-response}}

\textbf{Five most threatening prior works.} These were identified from
abstracts; full-text comparison is pending.

\begin{enumerate}
\def\labelenumi{\arabic{enumi}.}
\tightlist
\item
  CALIBURN \citep{youssef2026caliburn}: regime-dependent CRC for alert
  thresholds under an FP budget.
\item
  FIRCE \citep{barrett2026firce}: ``intrusion response'' plus conformal
  evaluation, but the response is retraining.
\item
  CORA and Conformal Policy Control \citep{cora2026, cpc2026}:
  risk-controlled gating of actions.
\item
  Hammar \& Stadler
  \citep{hammar2022optimalstopping, hammar2024dynamic}: optimal
  intrusion-response policies with threshold structure (IEEE TNSM).
\item
  SliceSecure and the sinkhole-slice idea \citep{khan2023slicesecure}:
  per-slice attack impact.
\end{enumerate}

A sixth, RBSID \citep{rbsid2025}, is risk-calibrated Bayesian decisions
and must be read.

\begin{longtable}[]{@{}
  >{\raggedright\arraybackslash}p{(\columnwidth - 18\tabcolsep) * \real{0.1000}}
  >{\centering\arraybackslash}p{(\columnwidth - 18\tabcolsep) * \real{0.1000}}
  >{\centering\arraybackslash}p{(\columnwidth - 18\tabcolsep) * \real{0.1000}}
  >{\centering\arraybackslash}p{(\columnwidth - 18\tabcolsep) * \real{0.1000}}
  >{\centering\arraybackslash}p{(\columnwidth - 18\tabcolsep) * \real{0.1000}}
  >{\centering\arraybackslash}p{(\columnwidth - 18\tabcolsep) * \real{0.1000}}
  >{\centering\arraybackslash}p{(\columnwidth - 18\tabcolsep) * \real{0.1000}}
  >{\centering\arraybackslash}p{(\columnwidth - 18\tabcolsep) * \real{0.1000}}
  >{\centering\arraybackslash}p{(\columnwidth - 18\tabcolsep) * \real{0.1000}}
  >{\raggedright\arraybackslash}p{(\columnwidth - 18\tabcolsep) * \real{0.1000}}@{}}
\caption{Claim-by-claim novelty matrix for C1 (\ding{51} covered,
\textasciitilde{} partially, \ding{55} not covered, ? unknown without full text,
--- not applicable).}\tabularnewline
\toprule\noalign{}
\begin{minipage}[b]{\linewidth}\raggedright
Claim
\end{minipage} & \begin{minipage}[b]{\linewidth}\centering
CALIBURN
\end{minipage} & \begin{minipage}[b]{\linewidth}\centering
FIRCE
\end{minipage} & \begin{minipage}[b]{\linewidth}\centering
CORA/CPC
\end{minipage} & \begin{minipage}[b]{\linewidth}\centering
LTT
\end{minipage} & \begin{minipage}[b]{\linewidth}\centering
Hammar--Stadler
\end{minipage} & \begin{minipage}[b]{\linewidth}\centering
RRE
\end{minipage} & \begin{minipage}[b]{\linewidth}\centering
SliceSecure
\end{minipage} & \begin{minipage}[b]{\linewidth}\centering
Jaqen
\end{minipage} & \begin{minipage}[b]{\linewidth}\raggedright
Status of claim
\end{minipage} \\
\midrule\noalign{}
\endfirsthead
\toprule\noalign{}
\begin{minipage}[b]{\linewidth}\raggedright
Claim
\end{minipage} & \begin{minipage}[b]{\linewidth}\centering
CALIBURN
\end{minipage} & \begin{minipage}[b]{\linewidth}\centering
FIRCE
\end{minipage} & \begin{minipage}[b]{\linewidth}\centering
CORA/CPC
\end{minipage} & \begin{minipage}[b]{\linewidth}\centering
LTT
\end{minipage} & \begin{minipage}[b]{\linewidth}\centering
Hammar--Stadler
\end{minipage} & \begin{minipage}[b]{\linewidth}\centering
RRE
\end{minipage} & \begin{minipage}[b]{\linewidth}\centering
SliceSecure
\end{minipage} & \begin{minipage}[b]{\linewidth}\centering
Jaqen
\end{minipage} & \begin{minipage}[b]{\linewidth}\raggedright
Status of claim
\end{minipage} \\
\midrule\noalign{}
\endhead
\bottomrule\noalign{}
\endlastfoot
K1 Finite-sample distribution-free guarantee & \ding{51} (alerts) &
\textasciitilde{} & \ding{51} & \ding{51} & \ding{55} & \ding{55} & \ding{55} & \ding{55} & \textbf{Not novel alone} \\
K2 Guarantee on service-level consequence of actions & \ding{55}? & \ding{55} &
\textasciitilde{} (generic harm) & --- & \ding{55} (model-based cost) & \ding{55} & \ding{55} &
\ding{55} & Plausibly novel {[}Inference{]} \\
K3 Heterogeneous per-slice budgets & \ding{55}? & \ding{55} & \ding{55} & --- & \ding{55} & \ding{55} & \ding{55} & \ding{55} &
Plausibly novel; thin alone \\
K4 Multi-action, non-monotone policy calibration & \ding{55}? & \ding{55} &
\textasciitilde{} (CPC) & \ding{51} (tool) & \ding{51} (not distribution-free) & \ding{51} & \ding{55} &
\ding{55} & Novel only in combination with K2 \\
K5 Cross-slice coupling via counterfactual replay or twin, with error
bound & \ding{55} & \ding{55} & \ding{55} & \ding{55} & \textasciitilde{} (emulation) & \ding{55} & \ding{55} & \ding{55} &
\textbf{Core novelty candidate} {[}Hypothesis{]} \\
K6 Validation on a sliced 5G core with measured SLA KPIs & \ding{55} & \ding{55} (IoT) &
\ding{55} & \ding{55} & \ding{55} (IT) & \ding{55} & \textasciitilde{} (detection) & \ding{55} & Novel \emph{if}
the testbed exists \\
\end{longtable}

\textbf{Incremental-novelty risk.} If K2--K5 reduce to ``apply LTT to a
new loss'', reviewers will call the work an application paper. The
contribution must therefore include:

\begin{itemize}
\tightlist
\item
  \textbf{(i)} a problem formulation in which collateral risk is defined
  against an \emph{oracle-mitigation counterfactual}, so that it
  isolates harm due to detection and action error;
\item
  \textbf{(ii)} a twin-in-the-loop calibration result that propagates
  surrogate error into the guarantee, with an empirical bound on that
  error;
\item
  \textbf{(iii)} a quantified analysis of temporal dependence across
  episodes;
\item
  \textbf{(iv)} evidence that alert-level control (K1 alone)
  \emph{fails} to bound per-slice collateral risk, the key falsifiable
  claim.
\end{itemize}

\textbf{Strongest likely reviewer objections.}

\begin{enumerate}
\def\labelenumi{\arabic{enumi}.}
\tightlist
\item
  ``This is LTT applied to networking.''
\item
  ``The i.i.d.-episode assumption is unrealistic for traffic.''
\item
  ``Emulated RAN, so no 5G claims.''
\item
  ``The counterfactual replay is expensive; its fidelity is unproven.''
\item
  ``The baselines are straw men.''
\end{enumerate}

The plan in Section \protect\hyperlink{sec:rec}{8} answers each one.

\hypertarget{c3-sla-weighted-multi-slice-quickest-detection}{%
\subsection{C3: SLA-weighted multi-slice quickest
detection}\label{c3-sla-weighted-multi-slice-quickest-detection}}

\textbf{Threats.} Multistream QCD with second-order optimality
\citep{fellouris2016second}; adaptive CUSUM for intrusions
\citep{tartakovsky2006novel}; an unpublished multi-slice QCD simulation
reporting cross-slice coherence effects \citep{qcdslice_github}. Fewer
than five threats were found because QCD-for-5G-slicing is sparsely
published, but the mature general theory is itself the main threat.

\textbf{Novelty matrix (summary).} Heterogeneous false-alarm constraints
per stream: \emph{unknown}, possibly covered in the multichannel
literature (not searched in depth). SLA-weighted delay: plausibly novel.
Coupled streams due to shared resources: partially addressed in
\citep{qcdslice_github}.

\textbf{Objection.} ``It is a corollary of known multichannel CUSUM
results.'' \textbf{Verdict:} survives as the backup only.

\hypertarget{c6-calibrated-gnn-digital-twin}{%
\subsection{C6: calibrated GNN digital
twin}\label{c6-calibrated-gnn-digital-twin}}

\textbf{Threats.} RouteNet-Fermi \citep{ferriolgalmes2023routenetfermi}
and the NDT vision paper \citep{almasan2022ndt}. We have \textbf{not}
searched specifically for conformal or uncertainty-aware GNN network
models; this is a known gap, G-07. \textbf{Verdict:} survives
provisionally, ranked third.

\hypertarget{sec:feas}{%
\section{Stage E --- Real-world feasibility audit}\label{sec:feas}}

\begin{longtable}[]{@{}
  >{\raggedright\arraybackslash}p{(\columnwidth - 4\tabcolsep) * \real{0.3333}}
  >{\raggedright\arraybackslash}p{(\columnwidth - 4\tabcolsep) * \real{0.3333}}
  >{\raggedright\arraybackslash}p{(\columnwidth - 4\tabcolsep) * \real{0.3333}}@{}}
\caption{Feasibility audit.}\tabularnewline
\toprule\noalign{}
\begin{minipage}[b]{\linewidth}\raggedright
Resource
\end{minipage} & \begin{minipage}[b]{\linewidth}\raggedright
Status
\end{minipage} & \begin{minipage}[b]{\linewidth}\raggedright
Evidence
\end{minipage} \\
\midrule\noalign{}
\endfirsthead
\toprule\noalign{}
\begin{minipage}[b]{\linewidth}\raggedright
Resource
\end{minipage} & \begin{minipage}[b]{\linewidth}\raggedright
Status
\end{minipage} & \begin{minipage}[b]{\linewidth}\raggedright
Evidence
\end{minipage} \\
\midrule\noalign{}
\endhead
\bottomrule\noalign{}
\endlastfoot
Compute (this session) & 4 vCPU, 15 GB RAM, no GPU, no LaTeX toolchain &
\textbf{{[}Verified{]}} by inspection \\
Network egress (this session) & Search only; scholarly APIs and arXiv
blocked & \textbf{{[}Verified{]}} by proxy logs \\
Prior manuscripts and code & Repository empty &
\textbf{{[}Verified{]}} \\
5G-NIDD & Public (IEEE DataPort, DOI 10.21227/xtep-hv36) &
\textbf{{[}Verified{]}}, record only \\
5GC PFCP dataset; SliceSecure dataset & Public on DataPort &
\textbf{{[}Verified{]}}, record only; label-quality concern noted \\
5G-SliciNdd (team dataset) & Not provided & \textbf{{[}Unverified{]}} \\
Team 5G testbed (eMBB/URLLC/mMTC shaping) & Described in the brief, not
inspected & \textbf{{[}Unverified{]}}: core, RAN, UPF, orchestration,
isolation, and mitigation hooks all unknown \\
Open-source sliced-core scaffolding & Open5GS Docker with per-slice UPF
\citep{xu2025rogue} & \textbf{{[}Verified{]}} that the repository is
referenced \\
RAN / SDR / RIC & None found & \textbf{{[}Unverified{]}}; treated as
unavailable \\
\end{longtable}

\textbf{Choice of validation method for C1.} The claims concern
\emph{core- and user-plane mitigation and the resulting per-slice KPIs}.
An \textbf{emulated sliced 5G core} (Open5GS or free5GC with UERANSIM, a
per-slice UPF, and tc/nftables enforcement) can test them directly.
Radio realism is not required, provided the claims are restricted to
that scope, and the emulated RAN will be stated as a threat to validity.
It is combined with \textbf{trace-driven replay} from 5G-NIDD and the
PFCP dataset for scale and repeated splits. Pure simulation is
insufficient, because the collateral loss depends on real queueing and
enforcement behaviour. A physical RAN is \emph{not} required to test the
hypotheses. \textbf{No new equipment is essential.} Access to and an
audit of the team testbed is the main dependency (G-01).

\hypertarget{sec:rec}{%
\section{Stage F --- Provisional recommendation}\label{sec:rec}}

\hypertarget{primary-provisional-c1}{%
\subsection{Primary (provisional): C1}\label{primary-provisional-c1}}

\begin{enumerate}
\def\labelenumi{\arabic{enumi}.}
\tightlist
\item
  \textbf{Title (working).} \emph{Certified Collateral-Risk Control for
  Slice-Aware Intrusion Response in 5G Networks}.
\item
  \textbf{Target journal.} IEEE Transactions on Network and Service
  Management (TNSM). Its scope covers security management, closed-loop
  automation, and service assurance, and it has published the closest
  response work \citep{hammar2022optimalstopping}. The alternative is
  IEEE TIFS if the theoretical component turns out to be stronger.
\item
  \textbf{Motivation and problem.} See C1 in Section
  \protect\hyperlink{sec:candidates}{4}.
\item
  \textbf{Closest work and remaining gap.} See Table 4.
\item
  \textbf{Formal hypothesis {[}Hypothesis{]}.}

  \begin{itemize}
  \tightlist
  \item
    \textbf{H1 (validity):} under exchangeable calibration episodes, the
    calibrated policy keeps the realised per-slice collateral risk at or
    below \(\alpha_s\) for every slice, with probability at least
    \(1-\delta\).
  \item
    \textbf{H2 (necessity):} a policy built on alert-level FPR control
    or pooled calibration, at the same nominal level, violates the URLLC
    collateral budget in a measurable fraction of splits.
  \item
    \textbf{H3 (utility):} at equal certified collateral budgets, the
    method's residual attack impact is no worse than that of a
    model-based threshold response policy
    \citep{hammar2022optimalstopping} tuned on the same data.
  \end{itemize}
\item
  \textbf{Candidate contributions.}

  \begin{itemize}
  \tightlist
  \item
    \textbf{(a)} The formulation: oracle-counterfactual collateral risk
    per slice.
  \item
    \textbf{(b)} Twin-in-the-loop LTT calibration with surrogate-error
    propagation.
  \item
    \textbf{(c)} A dependence- and shift-aware extension with a
    monitored recalibration trigger.
  \item
    \textbf{(d)} An open, reproducible sliced-core mitigation testbed
    and benchmark protocol.
  \end{itemize}
\item
  \textbf{Methodology.} Detector-agnostic: supervised gradient boosting
  and unsupervised Kitsune \citep{mirsky2018kitsune}. A structured
  threshold-policy family over actions. LTT with Hoeffding--Bentkus
  p-values and fixed-sequence or graph-based FWER control
  \citep{angelopoulos2025ltt}.
\item
  \textbf{Formulation (preliminary).} Let \(s\in\mathcal{S}\) index
  slices. Let \(e\) be an episode (a time window), \(\mathcal{B}_s(e)\)
  the set of benign UEs of slice \(s\), and \(\pi_\lambda\) a policy
  that maps detector scores and slice identity to actions in
  \(\mathcal{A}=\{\emptyset,\mathsf{RL}(\rho),\mathsf{Q},\mathsf{B}\}\).
  With \(V_u(e;\pi)\) the SLA-violation indicator of UE \(u\) under
  policy \(\pi\), and \(\pi^\star\) oracle mitigation, define
  \[L^{c}_s(e;\lambda)=\frac{1}{|\mathcal{B}_s(e)|}\sum_{u\in\mathcal{B}_s(e)}\mathbb{1}\{V_u(e;\pi_\lambda)=1,\;V_u(e;\pi^\star)=0\},\qquad R^{c}_s(\lambda)=\mathbb{E}_e[L^{c}_s(e;\lambda)].\]
  Let \(R^{a}(\lambda)\) be the expected residual attack impact. The
  target is
  \[\hat{\lambda}=\arg\min_{\lambda\in\hat{\Lambda}}\hat{R}^{a}(\lambda),\qquad \hat{\Lambda}=\{\lambda:\ H_{s,\lambda}:R^c_s(\lambda)>\alpha_s\ \mathrm{rejected}\ \forall s\},\]
  with FWER control at level \(\delta\) over all pairs \((s,\lambda)\),
  which yields
  \(\Pr\{\max_s(R^c_s(\hat{\lambda})-\alpha_s)\le 0\}\ge 1-\delta\).
  This follows from LTT \citep{angelopoulos2025ltt} and is \emph{not}
  claimed as novel. \textbf{Proposition to be proved {[}Hypothesis{]}:}
  if a twin \(\tilde{L}\) satisfies
  \(|\tilde{L}^c_s-L^c_s|\le\varepsilon_s\) with probability at least
  \(1-\eta\), then calibrating at \(\alpha_s-\varepsilon_s\) certifies
  the real system at \(\alpha_s\) with probability at least
  \(1-\delta-\eta\). \textbf{To be developed:} a degradation term for
  \(\beta\)-mixing episodes or for TV-bounded shift, following
  \citep{barber2023beyond}.
\item
  \textbf{Architecture.} Detector $\rightarrow$ score bus $\rightarrow$ policy engine
  (calibrated thresholds per slice and action) $\rightarrow$ enforcement at the UPF
  (N3/N6 nftables and tc per UE, keyed by TEID or IP; quarantine via
  re-association to a sinkhole S-NSSAI) $\rightarrow$ KPI telemetry $\rightarrow$ calibration
  store and drift monitor (for example OWAD-style \citep{han2023owad}).
\item
  \textbf{Data and testbed.} 5G-NIDD and PFCP traces for the replay arm.
  Team testbed or Open5GS-based emulation for counterfactual episodes.
  Benign traffic profiles per slice (eMBB bulk, URLLC periodic small
  packets, mMTC sparse).
\item
  \textbf{Baselines.}

  \begin{itemize}
  \tightlist
  \item
    B0: no mitigation.
  \item
    B1: fixed-threshold block.
  \item
    B2: conformal FPR-controlled detector with block (CALIBURN or
    Bates-style).
  \item
    B3: pooled, non-slice LTT.
  \item
    B4: cost-sensitive Bayes rule.
  \item
    B5: optimal-stopping threshold response (T-SPSA)
    \citep{hammar2022optimalstopping}.
  \item
    B6: oracle (upper bound).
  \end{itemize}

  \textbf{Ablations:} remove slice stratification; remove twin-error
  correction; single-action versus multi-action; vary calibration size
  \(n\). \textbf{Metrics:} realised per-slice collateral risk; frequency
  with which the guarantee is violated across at least 200 random and
  temporal splits; residual attack bytes; URLLC p99 latency; eMBB
  goodput; decision latency.
\item
  \textbf{Limitations and threats to validity.}

  \begin{itemize}
  \tightlist
  \item
    Emulated RAN: no radio claims.
  \item
    Exchangeability of episodes.
  \item
    Limited attack diversity (volumetric floods and scans; PFCP).
  \item
    Adaptive attackers who game the thresholds.
  \item
    Twin fidelity.
  \end{itemize}
\item
  \textbf{Feasibility.} Moderate. All software is open source; the main
  effort is the replay engine.
\item
  \textbf{Major risks and mitigations.}

  \begin{itemize}
  \tightlist
  \item
    The certified set may be empty at a realistic \(n\) $\rightarrow$ use
    fixed-sequence testing, coarser policy grids, or larger budgets for
    mMTC.
  \item
    Replay may be too slow $\rightarrow$ use a surrogate twin with an
    \(\varepsilon\) bound.
  \item
    A full read may show that a threat paper (CALIBURN, RBSID, CPC)
    already covers K2 $\rightarrow$ fall back to C3.
  \end{itemize}
\item
  \textbf{Minimum evidence before submission.}

  \begin{itemize}
  \tightlist
  \item
    H1 holds empirically on both replay and emulation.
  \item
    H2 is demonstrated, or the paper is re-scoped.
  \item
    The proposition is proved.
  \item
    Code and data are released.
  \item
    A full-text novelty audit against all six threats is complete.
  \end{itemize}
\end{enumerate}

\textbf{Why C1 over the alternatives.} It is the only candidate that

\begin{itemize}
\tightlist
\item
  \textbf{(i)} survives the hard-rejection criteria with an identifiable
  contribution,
\item
  \textbf{(ii)} reuses the team's assets (slice-stratified security,
  uncertainty, and the slice-differentiated testbed) without duplicating
  detection-level work, and
\item
  \textbf{(iii)} can be tested with infrastructure that is plausibly
  available.
\end{itemize}

\hypertarget{backup-provisional-c3}{%
\subsection{Backup (provisional): C3}\label{backup-provisional-c3}}

\emph{SLA-weighted multi-slice quickest attack detection under
heterogeneous false-alarm constraints.} \textbf{Target:} IEEE TNSM, or
IEEE Transactions on Signal and Information Processing over Networks if
the theory dominates. Activate it if the C1 full-text audit finds K2 or
K5 already covered.

\hypertarget{sec:roadmap}{%
\section{Stage G (outline) --- Phased roadmap}\label{sec:roadmap}}

The detailed roadmap is withheld until the direction is approved, as the
brief requires. Gate criteria for the first two phases:

\begin{itemize}
\tightlist
\item
  \textbf{Phase 1 (novelty validation).} Obtain and read the full texts
  of the six C1 threats, plus targeted searches G-02 to G-07.

  \begin{itemize}
  \tightlist
  \item
    \textbf{Go} if K2 and K5 remain uncovered.
  \item
    \textbf{Stop or switch to C3} otherwise.
  \end{itemize}
\item
  \textbf{Phase 2 (formulation).} Prove the twin-error proposition and
  specify the dependence model.

  \begin{itemize}
  \tightlist
  \item
    \textbf{Go} if the proof holds and the certified set is non-empty in
    a pilot replay of 5G-NIDD with n of 500 episodes or fewer.
  \end{itemize}
\item
  \textbf{Phases 3--8} follow the brief: baselines, implementation,
  experiments, ablation, writing, and audit.
\end{itemize}

Effort estimate (assumes one full-time researcher with access to the
testbed): Phases 1--2 in 4--6 weeks; Phases 3--6 in 4--6 months.
\textbf{{[}Hypothesis{]}}

\hypertarget{sec:gaps}{%
\section{Evidence and access gaps}\label{sec:gaps}}

\begin{longtable}[]{@{}
  >{\raggedright\arraybackslash}p{(\columnwidth - 6\tabcolsep) * \real{0.2500}}
  >{\raggedright\arraybackslash}p{(\columnwidth - 6\tabcolsep) * \real{0.2500}}
  >{\raggedright\arraybackslash}p{(\columnwidth - 6\tabcolsep) * \real{0.2500}}
  >{\raggedright\arraybackslash}p{(\columnwidth - 6\tabcolsep) * \real{0.2500}}@{}}
\caption{Evidence and access gaps.}\tabularnewline
\toprule\noalign{}
\begin{minipage}[b]{\linewidth}\raggedright
ID
\end{minipage} & \begin{minipage}[b]{\linewidth}\raggedright
Gap
\end{minipage} & \begin{minipage}[b]{\linewidth}\raggedright
Blocks
\end{minipage} & \begin{minipage}[b]{\linewidth}\raggedright
Action needed
\end{minipage} \\
\midrule\noalign{}
\endfirsthead
\toprule\noalign{}
\begin{minipage}[b]{\linewidth}\raggedright
ID
\end{minipage} & \begin{minipage}[b]{\linewidth}\raggedright
Gap
\end{minipage} & \begin{minipage}[b]{\linewidth}\raggedright
Blocks
\end{minipage} & \begin{minipage}[b]{\linewidth}\raggedright
Action needed
\end{minipage} \\
\midrule\noalign{}
\endhead
\bottomrule\noalign{}
\endlastfoot
G-01 & Team testbed not inspected (core, RAN, UPF, isolation, mitigation
hooks) & Stage E, C1 feasibility & Provide configs, topology, and logs,
or grant access \\
G-02 & Full text of CALIBURN (and its retitled version), RBSID, CPC,
CORA & C1 novelty (K2, K4) & Allow arxiv.org in the network policy, or
upload the PDFs \\
G-03 & Full text of Hammar \& Stadler TNSM 2022 / 2024 and TIFS 2025 &
C1 baselines and novelty & IEEE Xplore access \\
G-04 & CALM-5G (arXiv 2610.17962) existence and content & C1 / C2
novelty & Resolve the identifier \\
G-05 & DOIs for the 20 V-META and 11 PARTIAL entries & Bibliography
integrity & Allow doi.org and Crossref \\
G-06 & Team's prior manuscripts (5G-SliciNdd, leakage audit,
slice-stratified security) & Overlap rule & Upload them \\
G-07 & Targeted search on conformal or uncertainty-aware GNN network
twins; multichannel QCD with heterogeneous constraints & C6, C3 novelty
& Phase 1 \\
G-08 & Whether C-PP-COAD used 5G-NIDD (sources conflict) & C2 rejection
evidence (non-critical) & Full text \\
\end{longtable}

\hypertarget{sec:conc}{%
\section{Conclusion}\label{sec:conc}}

Within the stated verification limits, the evidence supports
\emph{abandoning} detection-level conformal false-alarm control as a
primary topic. It supports pursuing \textbf{certified collateral-risk
control for slice-aware intrusion response (C1)} as a provisional
primary direction, with \textbf{C3} as the backup. The checkpoint stops
here, pending approval and the closure of gaps G-01 to G-03.


\bmhead{Acknowledgements}
Literature screening, structuring and drafting were carried out with the assistance of an AI research
assistant (Claude Code). All bibliographic records carry an explicit verification status; no reference,
DOI or result was generated without a retrieved source.

\section*{Declarations}
\begin{itemize}
\item \textbf{Funding:} Not applicable / to be completed by the authors.
\item \textbf{Competing interests:} The authors declare no competing interests (to be confirmed).
\item \textbf{Ethics approval and consent to participate:} Not applicable (no human or animal data).
\item \textbf{Data availability:} The structured bibliography, literature matrix and scoring data are in
  \texttt{research/checkpoint1/data/}. No experimental data were generated at this checkpoint.
\item \textbf{Code availability:} \texttt{research/checkpoint1/scripts/} (bibliography generation, ranking,
  report build).
\item \textbf{Author contribution:} To be completed by the authors.
\end{itemize}

\bibliography{references}

\end{document}
